\section{Results}
\label{sec:results}

Our experiments test whether heavy-tailed exponent computation extends to Vision Transformers with bounded variance. The results reveal a nuanced picture: the methodology works, but requires stratification controls.

\subsection{Main Results}

Table~\ref{tab:main} presents aggregate statistics for heavy-tailed exponent measurements across 53 ViT models.

\begin{table}[h]
\centering
\begin{tabular}{ll}
\toprule
\textbf{Metric} & \textbf{Value} \\
\midrule
Models processed & 53 \\
Mean $\alpha$ & 4.327 \\
Median $\alpha$ & 3.521 \\
$\sigma(\alpha)$ global & \textbf{2.868} \\
$\alpha$ range & [2.20, 18.92] \\
\bottomrule
\end{tabular}
\caption{Aggregate statistics for heavy-tailed exponent measurements.}
\label{tab:main}
\end{table}

\textbf{Gate verdict: FAIL.} The global variance $\sigma = 2.868$ exceeds our threshold of 0.5 by 5.7$\times$.

\subsection{Within-Family Variance}

When we stratify by model family (organizational training origin), variance drops dramatically.

\begin{table}[h]
\centering
\begin{tabular}{lcc}
\toprule
\textbf{Family} & \textbf{n} & $\sigma(\alpha)$ \\
\midrule
google/vit-* & 6 & \textbf{0.24} \\
Global (all models) & 53 & 2.868 \\
\bottomrule
\end{tabular}
\caption{Within-family variance analysis.}
\label{tab:family}
\end{table}

This represents an \textbf{up to 12$\times$ variance reduction} through family stratification (based on google/vit-* family, n=6).

\subsection{Outlier Analysis}

Several models exhibit extreme $\alpha > 10$:

\begin{table}[h]
\centering
\begin{tabular}{llc}
\toprule
\textbf{Model} & \textbf{Training Task} & $\alpha$ \\
\midrule
jaranohaal/vit-base-violence-detection & Violence detection & 12.6 \\
mlx-vision/vit\_base\_patch16\_224-mlxim & Unknown fine-tuning & 14.8 \\
\bottomrule
\end{tabular}
\caption{Outlier models with extreme $\alpha$ values.}
\label{tab:outliers}
\end{table}

\textbf{Interpretation.} Outliers correspond to task-specific fine-tuning, not architectural variants.

\subsection{Summary of Findings}

\begin{table}[h]
\centering
\begin{tabular}{p{4cm}p{2.5cm}p{5cm}}
\toprule
\textbf{Research Question} & \textbf{Answer} & \textbf{Evidence} \\
\midrule
RQ1: Can $\alpha$ be computed for ViT? & \textbf{Yes} & 53/100 models processed successfully \\
RQ2: Is variance bounded ($\sigma < 0.5$)? & \textbf{No} globally, \textbf{Yes} within families & $\sigma = 2.868$ global, $\sigma = 0.24$ family \\
RQ3: Does family control help? & \textbf{Yes} & 12$\times$ variance reduction \\
\bottomrule
\end{tabular}
\caption{Summary of findings by research question.}
\label{tab:summary}
\end{table}
